\section{Waiting for fulfilment, receiving conversion}

The two interpretations answer different questions about the same
benefits. The first distinguishes deliverance already received from
the final gathering and resurrection still awaited. Augustine's
prophetic tense and Chrysostom's account of the Gospel explain why
continued danger need not contradict God's promise. The second
distinguishes confident petition from the demand that an unchanged
will receive whatever it wants. Aquinas explains that God gives the
faithful the capacity to live worthily. Theophylact describes the
petitioner's trust and Bede explains how that trust takes effect
when someone seeking pardon forgives.

Their practical consequences differ. The first concerns endurance
and discernment when the promised good is delayed or a false
claimant offers immediate assurance. The second concerns gratitude,
forgiveness and the willingness to receive the cure one requests.
An afflicted person needs more than an instruction to change
desires; a resentful petitioner needs more than an assurance
that fulfilment will eventually come.

Bellarmine's remembered rescue remains a genuine alternative
to Augustine's explanation of the Gradual's perfect tense.
That exegetical difference does not divide the two larger
interpretations into opposed doctrines. Both depend on
Christ's redemption and end in the inheritance for which
the faithful have been made fit. The Mass concludes with
people who have received asking that this gift continue
to heal them.
